\section{Анализ предметной области}

\subsection{Характеристика деятельности учебно-методического центра повышения квалификации}

Учебно-методические центры повышения квалификации и профессиональной переподготовки играют ключевую роль в обеспечении профессионального развития сотрудников различных отраслей. В условиях стремительных изменений в технологиях, рыночных реалиях необходимость в непрерывном обучении становится особенно важной. Профессионалы, чья деятельность напрямую зависит от актуальности знаний, обязаны поддерживать высокий уровень своей квалификации, чтобы оставаться конкурентоспособными и соответствовать современным требованиям.

Образовательные центры, предлагающие курсы повышения квалификации и переподготовки, становятся важнейшими инструментами для поддержания профессиональной пригодности специалистов в различных областях. Такие центры организуют обучение по дополнительным профессиональным программам, ведут учёт слушателей, преподавателей, учебных групп, а также формируют отчётность.

\subsection{Проблематика традиционного подхода к организации учебного процесса}

В традиционном подходе к организации учебных курсов и образовательных программ существует ряд значительных проблем, которые существенно снижают эффективность работы учебно-методического центра.

\textbf{Отсутствие централизованного хранилища данных.} В настоящее время информация о слушателях, преподавателях, программах повышения квалификации и профессиональной переподготовки хранится в разрозненных электронных таблицах. Это приводит к дублированию данных, затрудняет их поиск и актуализацию. Например, при изменении статуса слушателя или выдаче документа необходимо вручную обновлять несколько файлов, что повышает риск возникновения ошибок.

\textbf{Трудоёмкость ведения отчётности.} Учебно-методический центр обязан регулярно предоставлять статистическую отчётность. Для формирования такой отчётности требуется вручную агрегировать данные из множества источников: списков слушателей, журналов учёта документов, финансовых ведомостей. Ручной подсчёт показателей (число слушателей по категориям, источникам финансирования, возрасту, полу) занимает значительное время и чреват арифметическими ошибками.

\textbf{Сложность контроля за успеваемостью и статусом слушателей.} Отсутствие единой информационной системы не позволяет оперативно отслеживать, какие программы освоены слушателем, какие документы ему выданы, является ли он педагогическим работником или руководителем. Такая информация необходима как для внутреннего учёта, так и для ответов на запросы контролирующих органов. В текущих условиях получение подобных сведений требует ручного просмотра нескольких таблиц и сопоставления данных.

\textbf{Низкая скорость обработки запросов.} При обращении сотрудника центра с запросом типа «найти всех слушателей по фамилии Иванов, обучавшихся в 2024 году» или «выдать список слушателей с определённым статусом» поиск выполняется вручную с использованием штатных средств электронных таблиц. При большом объёме данных такой подход становится крайне неэффективным и может занимать много времени.

\textbf{Отсутствие разграничения прав доступа.} Хранение данных в общих файлах на сетевом диске не позволяет гибко настраивать права доступа: администратор, оператор и простой пользователь имеют одинаковые возможности по просмотру и редактированию информации. Это создаёт риск случайного или преднамеренного изменения важных данных.

Таким образом, существующий «ручной» подход к управлению учебным процессом исчерпал свои возможности. Дальнейшее развитие центра требует внедрения специализированной информационной системы, которая автоматизирует учётные операции, ускорит формирование отчётности и повысит достоверность данных.

\subsection{Обоснование необходимости автоматизации}

Перечисленные проблемы могут быть эффективно решены путём создания информационной системы для автоматизации процессов управления учебными курсами. Такая система должна обеспечить:

\begin{itemize}
	\item централизованное хранение и управление данными о слушателях, преподавателях, программах и группах;
	\item возможность гибкой фильтрации и поиска информации по множеству критериев;
	\item автоматическое формирование статистических сводок;
	\item импорт данных из электронных таблиц (например, из файлов Excel, предоставляемых заказчиком);
	\item разграничение прав доступа (администраторы, операторы);
	\item удобный веб-интерфейс для работы с системой.
\end{itemize}

Автоматизация позволит централизовать все данные, устранить ошибки при ручном вводе, обеспечить оперативный контроль за успеваемостью слушателей и существенно сократить время на подготовку отчётности.

\subsection{Анализ существующих систем автоматизации учебного процесса}

На рынке представлен широкий спектр информационных систем для управления обучением. Ниже рассмотрены наиболее распространённые решения, их функциональные возможности и ограничения применительно к задачам учебно-методического центра повышения квалификации.

\subsubsection{Moodle}
Moodle является одной из самых популярных бесплатных систем управления обучением с открытым исходным кодом. Основные функции Moodle включают:
\begin{itemize}
	\item создание и управление курсами, лекциями, тестами и заданиями;
	\item форум, чат, обмен сообщениями между преподавателями и студентами;
	\item ведение журнала оценок и отслеживание активности пользователей;
	\item поддержка SCORM и других стандартов дистанционного обучения;
	\item разграничение ролей (администратор, преподаватель, студент).
\end{itemize}

\textbf{Избыточность для задач центра.} Moodle ориентирован на организацию полноценного дистанционного обучения с интерактивными элементами. Для учебно-методического центра, основная задача которого – учёт слушателей, программ и формирование отчётности, функционал Moodle является чрезмерным. Кроме того, Moodle не предоставляет готовых инструментов для формирования статистической отчётности, а его настройка под специфические требования потребует значительных временных и технических ресурсов.

\subsubsection{Google Classroom}
Google Classroom – это бесплатный сервис от Google, интегрированный с Google Диском, Gmail и Google Календарём. Основные возможности:
\begin{itemize}
	\item создание классов и заданий, прикрепление материалов из Google Диска;
	\item автоматическое создание папок для каждого класса и студента;
	\item обратная связь и оценки;
	\item интеграция с другими сервисами Google.
\end{itemize}

\textbf{Избыточность для задач центра.} Google Classroom не предоставляет гибких средств для ведения детального учёта слушателей. Отсутствуют механизмы формирования сложных отчётов и импорта данных из внешних систем. Сервис не предназначен для работы с персональными данными в объёме, требуемом для отчётности перед государственными органами.

\subsubsection{1С:Университет}
Программный продукт «1С:Университет» разработан для автоматизации деятельности образовательных учреждений высшего и среднего профессионального образования. Его функционал включает:
\begin{itemize}
	\item учёт контингента студентов, успеваемости, посещаемости;
	\item формирование учебных планов и рабочих программ;
	\item расчёт нагрузки преподавателей;
	\item подготовка государственных отчётов (в том числе для Росстата);
	\item интеграция с системами «1С:Зарплата и управление персоналом».
\end{itemize}

\textbf{Избыточность для задач центра.} Система обладает высокой стоимостью внедрения и сопровождения. Она избыточна для небольшого учебно-методического центра, поскольку содержит множество модулей, не имеющих отношения к дополнительному профессиональному образованию. Кроме того, «1С:Университет» требует установки и администрирования специализированного серверного программного обеспечения, что повышает порог входа.

\subsubsection{Вывод по анализу аналогов}

Проведённый анализ существующих систем управления обучением показал, что ни одна из них в полной мере не соответствует специфике деятельности и организационно-техническим требованиям учебно-методического центра.

Универсальные образовательные платформы (Moodle, Google Classroom) обладают избыточным функционалом для организации дистанционного взаимодействия, однако не предоставляют специализированных средств для детального учёта контингента слушателей, ведения документов об образовании и формирования утверждённых форм статистической отчётности. Их адаптация под данные задачи потребовала бы значительной доработки, что сопоставимо с разработкой нового решения.

С другой стороны, комплексные системы, такие как «1С:Университет», предлагают широкие возможности для автоматизации деятельности образовательных учреждений, однако их внедрение сопряжено с высокой стоимостью программного обеспечения и необходимостью обслуживания квалифицированным персоналом. Функциональные возможности этих систем значительно превышают реальные потребности центра, что приводит к неоправданным затратам.

Таким образом, целесообразным является не внедрение готового продукта, а разработка собственной специализированной информационной системы. Это позволит создать решение, максимально адаптированное под уникальные бизнес-процессы центра, избежав при этом функциональной избыточности и высоких лицензионных расходов.

\subsection{Постановка задач на разработку}

На основе проведённого анализа были сформулированы следующие задачи, которые должна решать разрабатываемая информационная система:

\begin{enumerate}
	\item Обеспечить ввод, хранение и редактирование информации о слушателях, преподавателях, программах повышения квалификации и профессиональной переподготовки.
	\item Реализовать механизмы фильтрации и поиска данных по различным критериям.
	\item Предоставить возможность выгрузки статистических сводок.
	\item Реализовать импорт исходных данных из Excel-файлов (в соответствии с предоставленными образцами).
	\item Обеспечить разграничение прав доступа.
	\item Создать удобный веб-интерфейс для работы с системой и десктопное приложение для управления серверной частью.
\end{enumerate}